% Appendix: Luttinger–Tisza Method
% Converted on 2026-10-08 from docs/lswt/04-appendices/luttinger-tisza-method.md (status: draft, source-section: "None in the primary note; the restructured TeX contains only an empty TODO section (app:luttinger-tisza). Content follows Luttinger and Tisza (1946) and Lyons and Kaplan (1960).").
% Review status: draft; awaiting the user's physics and mathematics review.

\hypertarget{sec-lswt-luttinger-tisza-method}{%
\section{Appendix: Luttinger--Tisza Method}\label{sec-lswt-luttinger-tisza-method}}

LSWT expands about a classical reference configuration, and the expansion is meaningful only when that configuration minimizes the classical energy. The Luttinger--Tisza (LT) method gives a lower bound on the classical energy of a bilinear spin model by relaxing the fixed length of each spin to a single global constraint. The relaxed problem reduces to the diagonalization of a matrix in momentum space. When a configuration built from the minimizing eigenvectors also has unit spin length at every site, it attains the bound and is a classical ground state. This appendix derives the bound, the conditions for attaining it with a single wave vector, the generalization to inequivalent sublattices, and the limits of the method.

\subsection{Classical Energy as a Quadratic Form}

We consider the bilinear exchange and single-ion terms of \lswtnoteref{LSWT spin Hamiltonian}{Bilinear Spin Hamiltonian} without the Zeeman term. The classical energy is the expectation value of these terms in the spin-coherent product state with unit vectors \(\mathbf n_I\), as defined in \lswtnoteref{LSWT spin Hamiltonian}{LSWT Overview}. For the exchange terms this replaces each spin operator by the classical vector \(S_I\mathbf n_I\). For the single-ion term the coherent-state value is \(\langle\hat S_I^\alpha\hat S_I^\beta+\hat S_I^\beta\hat S_I^\alpha\rangle/2=S_I(S_I-\tfrac12)n_I^\alpha n_I^\beta+\tfrac{S_I}2\delta^{\alpha\beta}\), which differs from \(S_I^2n_I^\alpha n_I^\beta\) at finite spin and reduces to the constant \(\operatorname{Tr}\mathsf D_I/4\) at \(S_I=\tfrac12\). Writing the link sum as an ordered-pair sum gives

\[
E_{\mathrm{cl}}[\{\mathbf n_I\}]
=\frac12\sum_{I\neq J}S_IS_J\,\mathbf n_I^{\mathsf T}\mathsf J_{IJ}\mathbf n_J
+\sum_IS_I\big(S_I-\tfrac12\big)\,\mathbf n_I^{\mathsf T}\mathsf D_I\mathbf n_I
+E_D,
\qquad
E_D=\sum_I\frac{S_I}{2}\operatorname{Tr}\mathsf D_I,
\]

where \(\mathsf J_{JI}=\mathsf J_{IJ}^{\mathsf T}\) and the constant \(E_D\) does not depend on the spin directions. Apart from \(E_D\), the energy is a quadratic form in the \(3N_{\mathrm{site}}\) components of the unit vectors. The ground-state problem minimizes this form under the strong constraint

\[
|\mathbf n_I|=1\quad\text{for every site }I.
\]

The LT method replaces the \(N_{\mathrm{site}}\) strong constraints by the single weak constraint

\[
\sum_I|\mathbf n_I|^2=N_{\mathrm{site}},
\]

which every configuration satisfying the strong constraint also satisfies.

\subsection{Momentum-Space Form}

The LT problem uses the unit cell of the Hamiltonian, with sites \(I=(i,a)\) at \(\mathbf r_{ia}=\mathbf R_i+\boldsymbol\delta_a\) and \(a=1,\ldots,N_s\); the magnetic unit cell of the ground state is not known in advance. With the full-position Fourier convention of \lswtnoteref{LSWT boson Hamiltonian}{Momentum-Space BdG Hamiltonian},

\[
\mathbf n_{ia}=\frac{1}{\sqrt{N_{\mathrm{uc}}}}\sum_{\mathbf q}\exp(\mathrm i\mathbf q\cdot\mathbf r_{ia})\,\mathbf n_a(\mathbf q),
\qquad
\mathbf n_a(-\mathbf q)=\mathbf n_a(\mathbf q)^*,
\]

where the reality condition follows from real \(\mathbf n_{ia}\) and the momenta run over the Brillouin zone of the Hamiltonian's unit cell. Translation invariance makes the energy diagonal in \(\mathbf q\):

\begin{equation}\label{eq-lswt-lt-matrix}
E_{\mathrm{cl}}=\sum_{\mathbf q}\mathbf n(\mathbf q)^\dagger\,\mathsf L_{\mathbf q}\,\mathbf n(\mathbf q)+E_D,
\qquad
\big[\mathsf L_{\mathbf q}\big]_{ab}
=\frac12\sum_{J\in b}S_aS_b\,\mathsf J_{IJ}\exp\!\big(\mathrm i\mathbf q\cdot(\mathbf r_J-\mathbf r_I)\big)
+\delta_{ab}\,S_a\big(S_a-\tfrac12\big)\,\mathsf D_a,
\end{equation}

where \(\mathbf n(\mathbf q)\) collects the \(3N_s\) components \(\mathbf n_a(\mathbf q)\), \(I\) is a fixed site of sublattice \(a\), and the sum runs over the sites \(J\neq I\) of sublattice \(b\) coupled to \(I\). Each \([\mathsf L_{\mathbf q}]_{ab}\) is a \(3\times3\) block. The relation \(\mathsf J_{JI}=\mathsf J_{IJ}^{\mathsf T}\) makes the \(3N_s\times3N_s\) LT matrix \(\mathsf L_{\mathbf q}\) Hermitian, and real interaction matrices give \(\mathsf L_{-\mathbf q}=\mathsf L_{\mathbf q}^*\). The spin lengths are absorbed into \(\mathsf L_{\mathbf q}\), so the constraints refer to unit vectors for any \(S_a\). The weak constraint becomes \(\sum_{\mathbf q}\|\mathbf n(\mathbf q)\|^2=N_{\mathrm{site}}\).

\subsection{Luttinger--Tisza Lower Bound}

Let \(\lambda_{\min}(\mathbf q)\) be the lowest eigenvalue of \(\mathsf L_{\mathbf q}\) and \(\lambda_{\mathrm{LT}}=\min_{\mathbf q}\lambda_{\min}(\mathbf q)\). Each momentum component satisfies \(\mathbf n(\mathbf q)^\dagger\mathsf L_{\mathbf q}\mathbf n(\mathbf q)\ge\lambda_{\mathrm{LT}}\|\mathbf n(\mathbf q)\|^2\), so every configuration obeying the weak constraint, and in particular every configuration of unit spins, satisfies

\begin{equation}\label{eq-lswt-lt-bound}
\frac{E_{\mathrm{cl}}-E_D}{N_{\mathrm{site}}}\ge\lambda_{\mathrm{LT}}.
\end{equation}

Equality holds exactly when \(\mathbf n(\mathbf q)\) vanishes except at the minimizing wave vectors \(\mathbf Q\) and their negatives, and lies in the eigenspace of \(\lambda_{\mathrm{LT}}\) there. A configuration of this form that also satisfies the strong constraint attains the bound and is therefore a classical ground state, with energy \(N_{\mathrm{site}}\lambda_{\mathrm{LT}}+E_D\). If no such configuration exists, \(\lambda_{\mathrm{LT}}\) is a strict lower bound and the method does not determine the ground state.

\subsection{Single-Wave-Vector Configurations}

The simplest candidates use one minimizing wave vector \(\mathbf Q\) with an eigenvector \(\mathbf w=(\mathbf w_1,\ldots,\mathbf w_{N_s})\) of \(\lambda_{\mathrm{LT}}\). The component at \(-\mathbf Q\) is fixed by the reality condition, since \(\mathbf w^*\) is an eigenvector of \(\mathsf L_{-\mathbf Q}\). The configuration is

\[
\mathbf n_{ia}=\operatorname{Re}\!\big[\mathbf u_a\exp(\mathrm i\mathbf Q\cdot\mathbf R_i)\big],
\qquad
\mathbf u_a=\mathbf w_a\exp(\mathrm i\mathbf Q\cdot\boldsymbol\delta_a),
\]

with the cell amplitude \(\mathbf u_a\) and an overall scale of \(\mathbf w\) still free. Its squared length is

\begin{equation}\label{eq-lswt-lt-single-q-length}
|\mathbf n_{ia}|^2=\frac12|\mathbf u_a|^2+\frac12\operatorname{Re}\!\big[(\mathbf u_a\cdot\mathbf u_a)\exp(2\mathrm i\mathbf Q\cdot\mathbf R_i)\big],
\end{equation}

where \(\mathbf u_a\cdot\mathbf u_a=\sum_\gamma(u_a^\gamma)^2\) carries no complex conjugation. The phases \(\exp(2\mathrm i\mathbf Q\cdot\mathbf R_i)\) over all cells form a subgroup of the unit circle, and the strong constraint takes one of three forms:

\begin{itemize}
\tightlist
\item
  If \(2\mathbf Q\) is a reciprocal-lattice vector, as at the zone center or at half reciprocal vectors, then \(\exp(\mathrm i\mathbf Q\cdot\mathbf R_i)=\pm1\), the configuration is collinear within each sublattice, and the condition is \(|\operatorname{Re}\mathbf u_a|=1\) after a suitable choice of the overall phase of \(\mathbf w\).
\item
  If \(4\mathbf Q\) is a reciprocal-lattice vector but \(2\mathbf Q\) is not, the phases \(\exp(2\mathrm i\mathbf Q\cdot\mathbf R_i)\) take only the values \(\pm1\). The condition is \(\operatorname{Re}(\mathbf u_a\cdot\mathbf u_a)=0\) and \(|\mathbf u_a|^2=2\). Writing \(\mathbf u_a=\mathbf x_a+\mathrm i\mathbf y_a\), this requires \(|\mathbf x_a|=|\mathbf y_a|=1\) but leaves the angle between \(\mathbf x_a\) and \(\mathbf y_a\) free. The configuration cycles through \(\mathbf x_a,-\mathbf y_a,-\mathbf x_a,\mathbf y_a\); for \(\mathbf y_a=\pm\mathbf x_a\) it is the collinear up-up-down-down pattern.
\item
  Otherwise the phases take at least three values that are not confined to \(\pm1\), and the condition is \(\mathbf u_a\cdot\mathbf u_a=0\) and \(|\mathbf u_a|^2=2\). Then \(\mathbf x_a\) and \(\mathbf y_a\) are orthonormal, and each sublattice forms a planar spiral in the plane spanned by \(\mathbf x_a\) and \(\mathbf y_a\).
\end{itemize}

When the eigenspace of \(\lambda_{\mathrm{LT}}\) at \(\mathbf Q\) has more than one dimension, \(\mathbf w\) may be any vector in it, and the conditions become equations for its coefficients. For a Bravais lattice (\(N_s=1\)) with isotropic Heisenberg exchange, \(\mathsf L_{\mathbf q}=S^2J(\mathbf q)\,\mathsf I_3\) with \(J(\mathbf q)=\frac12\sum_{\boldsymbol\Delta}J_{\boldsymbol\Delta}\exp(\mathrm i\mathbf q\cdot\boldsymbol\Delta)\), where \(\boldsymbol\Delta\) runs over all neighbor vectors. Every complex vector is then an eigenvector, a planar spiral always satisfies the strong constraint, and a spiral with wave vector \(\mathbf Q\) is a ground state, as shown by Lyons and Kaplan. When several inequivalent wave vectors minimize \(J(\mathbf q)\), other ground states, including multiple-wave-vector states, can have the same energy. For the triangular-lattice antiferromagnet with nearest-neighbor exchange \(J>0\), \(J(\mathbf q)=J\sum_{m=1}^{3}\cos(\mathbf q\cdot\mathbf a_m)\) over the three bond directions has its minimum \(-\frac32J\) at the zone corner, and the spiral at that wave vector is the \(120^\circ\) state with energy \(-\frac32JS^2\) per site.

\subsection{Inequivalent Sublattices and the Generalized Method}

For \(N_s>1\), the minimizing eigenvector generally distributes its weight unequally among the sublattices, \(|\mathbf w_a|\neq|\mathbf w_b|\), while the strong constraint requires equal lengths. The single weak constraint then allows configurations that no unit-spin state can reach, and the bound is often not attained. A stronger bound follows from one Lagrange multiplier \(\lambda_a\) per sublattice. For unit spins, \(\sum_I\lambda_{a(I)}(|\mathbf n_I|^2-1)=0\), so

\[
E_{\mathrm{cl}}
=\sum_{\mathbf q}\mathbf n(\mathbf q)^\dagger\big(\mathsf L_{\mathbf q}-\Lambda\big)\mathbf n(\mathbf q)
+N_{\mathrm{uc}}\sum_{a=1}^{N_s}\lambda_a,
\qquad
\Lambda=\operatorname{diag}(\lambda_1,\ldots,\lambda_{N_s})\otimes\mathsf I_3.
\]

If \(\mathsf L_{\mathbf q}-\Lambda\) is positive semidefinite at every \(\mathbf q\), the first term is nonnegative and \(E_{\mathrm{cl}}\ge N_{\mathrm{uc}}\sum_a\lambda_a\). The best bound of this form maximizes \(\sum_a\lambda_a\) under that condition. Equal multipliers \(\lambda_a=\lambda_{\mathrm{LT}}\) recover \eqref{eq-lswt-lt-bound}, so the generalized bound is never weaker. Lyons and Kaplan introduced this generalization through adjustable parameters in the weak constraint, which allows crystallographically inequivalent spins.

\subsection{Scope and Limitations}

\begin{itemize}
\tightlist
\item
  \textbf{Zeeman field.} The Zeeman term is linear in the spin vectors, so the energy is no longer a homogeneous quadratic form, and the bound above does not include it.
\item
  \textbf{Single-ion anisotropy.} An anisotropic \(\mathsf D_I\) makes the diagonal blocks of \(\mathsf L_{\mathbf q}\) distinguish spin directions, so the minimizing eigenvectors generally do not combine into unit spins; for example, easy-plane anisotropy with a spiral whose plane contains the hard axis. For such models the LT result is a lower bound and a list of candidate wave vectors rather than a solution.
\item
  \textbf{Multiple wave vectors.} When no single-wave-vector configuration satisfies the strong constraint, a combination of several minimizing wave vectors may still do so. Failure of the single-wave-vector construction therefore does not show that the bound is unattainable.
\item
  \textbf{Degenerate minima.} When the bound is attained, a minimum of \(\lambda_{\min}(\mathbf q)\) on a line or an extended region, or a degenerate eigenspace beyond that required by a global spin-rotation symmetry, signals a degenerate manifold of classical ground states, and the choice among them is not made by the classical energy. The threefold degeneracy of \(\mathsf L_{\mathbf q}\) for isotropic exchange reflects only the global rotation of all spins. When the bound is not attained, the degeneracy of the LT minimum does not determine the degeneracy of the true ground states.
\end{itemize}

\subsection{Relation to Linear Spin-Wave Theory}

LSWT requires a reference configuration of unit spins that is stationary under the fixed-length constraint, as described in Classical Order and Local Frame (Section~\ref{sec-lswt-classical-order-and-local-frame}). A configuration that attains the LT bound is a global minimum of the zero-field classical energy. Every harmonic fluctuation about it therefore raises or preserves the energy, and its BdG matrix \(\mathsf H_{\mathbf k}\) is at least positive semidefinite, provided the quadratic boson Hamiltonian uses the same coherent-state value of the single-ion term as the classical energy above (\lswtnoteref{LSWT boson Hamiltonian}{Real-Space Boson Hamiltonian}), as required in \lswtnoteref{LSWT paraunitary diagonalization}{Paraunitary Diagonalization}. When the bound is not attained, \(\lambda_{\mathrm{LT}}\) still bounds the energy of any candidate configuration obtained by numerical minimization from below, and the minimizing wave vectors indicate the magnetic unit cells to consider.

\subsection{References}

\subsubsection{Internal Documents}

\begin{itemize}
\tightlist
\item
  \lswtnoteref{LSWT spin Hamiltonian}{Bilinear Spin Hamiltonian}: defines the exchange and single-ion matrices and the link counting used in the classical energy.
\item
  Classical Order and Local Frame (Section~\ref{sec-lswt-classical-order-and-local-frame}): defines the classical reference configuration.
\item
  \lswtnoteref{LSWT boson Hamiltonian}{Momentum-Space BdG Hamiltonian}: defines the full-position Fourier convention.
\item
  \lswtnoteref{LSWT paraunitary diagonalization}{Paraunitary Diagonalization}: states the positive-definiteness requirement on the BdG matrix.
\end{itemize}

\subsubsection{External Sources}

\begin{itemize}
\tightlist
\item
  J. M. Luttinger and L. Tisza, ``Theory of Dipole Interaction in Crystals,'' \emph{Physical Review} \textbf{70}, 954--964 (1946), \href{https://doi.org/10.1103/PhysRev.70.954}{doi:10.1103/PhysRev.70.954}: representation of spin arrays as vectors in a many-dimensional space and reduction of the minimum-energy problem to a matrix diagonalization.
\item
  D. H. Lyons and T. A. Kaplan, ``Method for Determining Ground-State Spin Configurations,'' \emph{Physical Review} \textbf{120}, 1580 (1960), \href{https://doi.org/10.1103/PhysRev.120.1580}{doi:10.1103/PhysRev.120.1580}: generalization with adjustable parameters in the weak constraint for inequivalent spins; spiral ground state for lattices of equivalent spins.
\end{itemize}
